\subsection{Removing a fixed additive error by one finite multiplier}

A frequent final step in a rational estimate is to remove an additive
rounding allowance $J$ that is independent of precision. Suppose the
discrepancy has already been cleared of denominators and its positive part
is $z\in\mathbb N$. The following are interderivable inside first-order
iPRA, for each $J$:
\[
 z=0,\qquad (J+1)z\le J,\qquad
                 \forall k\ k z\le J .
\]
One can even allow $J$ to be an existential witness. Thus
$z=0\leftrightarrow\exists J\,\forall k\,kz\le J$. The order of the
quantifiers matters. For $z=1$ the variant
$\forall k\,\exists J\,kz\le J$ is satisfied by $J=k$.
In the motivating comparison $ka\le kb+J$ for naturals $a,b$,
ordinary finite arithmetic rewrites the discrepancy as
$z=a\mathbin{\dot-} b$. Its positive part obeys $kz\le J$, hence $a\le b$.
The Lean derivation below treats the excess $z$ as its input.

\paragraph{An object proof of the separating multiplier.}
Use the countdown $D(j,x)$ from the preceding argument. Since
$D(j,x)=x\mathbin{\dot-}j$, the order test $x\le j$ is the quantifier-free
equation $D(j,x)=0$. Define $M(0,k)=0$ and
$M(z+1,k)=k+M(z,k)$. This computes $kz$ with suitably ordered summands.
First prove
\[
              D(j,(j+1)+y)=y+1
\]
by quantifier-free induction on $j$. At $j=0$ it is the defining
zero-step equation. For the successor, rewrite the input as
$1+((j+1)+y)$ and apply $D(j+1,x+1)=D(j,x)$.
It follows that at a positive excess,
\[
           D(J,M(z+1,J+1))=M(z,J+1)+1.
\]
The right side is a successor, so it cannot equal zero. For $z=0$ all
the multiples are zero. This proves the finite test at $J+1$ and, by
specializing a universal hypothesis, the universal test. An existential
budget is eliminated by this same derivation with the supplied $J$.
In the forward direction zero itself can serve as a budget.
Every induction in this proof has a quantifier-free invariant.

\texttt{PrecisionBounds.finite\_recovery},
\texttt{fixed\_budget} and \texttt{bounded\_iff} are object derivations.
Applying the last theorem to the excess in each of the three primary
sentences allows a separate budget for each packet and gives three more
closed biconditionals of type \texttt{Prf PRAAx []}. In particular,
no uniform budget across different number fields, conductors or cutoffs
is assumed. The multiplier $J+1$ eliminates the budget for its own packet.
